% from ega2-1-fr.tex
\oldpage[II]{159}

\begin{lemma}[8.2.7]
\label{II.8.2.7-fr}
Soient $S$ un anneau gradué à degrés positifs, $f$ un élément homogène de
$S_d$ ($d>0$). On a des isomorphismes canoniques d'anneaux
\begin{equation}
\label{II.8.2.7.1-fr}
  S_f\xrightarrow{\sim}\bigoplus_{n\in\mathbf{Z}}S(n)_{(f)}
\tag{8.2.7.1}
\end{equation}
\begin{equation}
\label{II.8.2.7.2-fr}
  (S_f^{\geq})_{f/1}\xrightarrow{\sim}S_f
\tag{8.2.7.2}
\end{equation}
\begin{equation}
\label{II.8.2.7.3-fr}
  S^\natural_{(f^\natural)}\xrightarrow{\sim}S_f^{\geq}
\tag{8.2.7.3}
\end{equation}
dont les deux premiers sont des isomorphismes d'anneaux gradués.
\end{lemma}

Il est immédiat par définition que l'on a
$(S_f)_n=(S(n)_f)_0$, d'où l'isomorphisme
(\hyperref[II.8.2.7.1-fr]{8.2.7.1}), qui n'est autre que l'identité. D'autre
part, comme $f/1$ est inversible dans $S_f$, il y a un isomorphisme canonique
$S_f\xrightarrow{\sim}(S_f^{\geq})_{f/1}=(S_f)_{f/1}$ en vertu de
(\hyperref[II.8.2.1.2-fr]{8.2.1.2}) appliqué à $S_f$~; l'isomorphisme
réciproque est par définition l'isomorphisme
(\hyperref[II.8.2.7.2-fr]{8.2.7.2}). Enfin, si
$x=\sum_{m\geq n}y_m$ est un élément de $S_{[n]}$, avec $n=kd$, on fait
correspondre à l'élément $x/(f^\natural)^k$ l'élément
$\sum_m y_m/f^k$ de $S_f^{\geq}$, et on vérifie aussitôt que cela définit un
isomorphisme (\hyperref[II.8.2.7.3-fr]{8.2.7.3}).

\begin{env}[8.2.8]
\label{II.8.2.8-fr}
Si $M$ est un $S$-module gradué, on pose de même pour tout $n\in\mathbf{Z}$
\begin{equation}
\label{II.8.2.8.1-fr}
  M_{[n]}=\bigoplus_{m\geq n}M_m
\tag{8.2.8.1}
\end{equation}
et comme $S_{[m]}M_{[n]}\subset M_{[m+n]}$ ($m\geq0$), on peut définir un
$S^\natural$-module gradué $M^\natural$ en prenant
\begin{equation}
\label{II.8.2.8.2-fr}
  M^\natural=\bigoplus_{n\in\mathbf{Z}}M_n^\natural,
  \quad\text{avec}\quad M_n^\natural=M_{[n]}.
\tag{8.2.8.2}
\end{equation}
\end{env}

Nous laissons encore au lecteur la démonstration du~:

\begin{lemma}[8.2.9]
\label{II.8.2.9-fr}
Avec les notations de (\hyperref[II.8.2.7-fr]{8.2.7}) et
(\hyperref[II.8.2.8-fr]{8.2.8}), on a les di-isomorphismes canoniques de
modules
\begin{equation}
\label{II.8.2.9.1-fr}
  M_f\xrightarrow{\sim}\bigoplus_{n\in\mathbf{Z}}M(n)_{(f)}
\tag{8.2.9.1}
\end{equation}
\begin{equation}
\label{II.8.2.9.2-fr}
  (M_f^{\geq})_{f/1}\xrightarrow{\sim}M_f
\tag{8.2.9.2}
\end{equation}
\begin{equation}
\label{II.8.2.9.3-fr}
  M^\natural_{(f^\natural)}\xrightarrow{\sim}M_f^{\geq}
\tag{8.2.9.3}
\end{equation}
dont les deux premiers sont des di-isomorphismes de modules gradués.
\end{lemma}

\begin{lemma}[8.2.10]
\label{II.8.2.10-fr}
Soit $S$ un anneau gradué à degrés positifs.
\begin{enumerate}
\item[\rm(i)] Pour que $S^\natural$ soit une $S_0^\natural$-algèbre de type
fini (resp. noethérienne) il faut et il suffit que $S$ soit une $S_0$-algèbre
de type fini (resp. noethérienne).
\item[\rm(ii)] Pour que $S_{n+1}^\natural=S_1^\natural S_n^\natural$ pour
$n\geq n_0$, il faut et il suffit que $S_{n+1}=S_1S_n$ pour $n\geq n_0$.
\item[\rm(iii)] Pour que $S_n^\natural=S_1^{\natural n}$ pour $n\geq n_0$, il
faut et il suffit que $S_n=S_1^n$ pour $n\geq n_0$.
\item[\rm(iv)] Si $(f_\alpha)$ est un ensemble d'éléments homogènes de $S_+$
telle que $S_+$ soit la racine dans $S_+$ de l'idéal de $S_+$ engendré par
les $f_\alpha$, alors $S_+^\natural$ est la racine dans $S_+^\natural$ de
l'idéal de $S_+^\natural$ engendré par les $f_\alpha^\natural$.
\end{enumerate}
\end{lemma}

\begin{enumerate}
\item[\rm(i)] Si $S^\natural$ est une $S_0^\natural$-algèbre de type fini,
$S_+=S_1^\natural$ est un module de type fini sur $S=S_0^\natural$ par
(\hyperref[II.2.1.6-fr]{2.1.6}, (i)), donc $S$ est une $S_0$-algèbre de type
fini (\hyperref[II.2.1.4-fr]{2.1.4})~; si $S^\natural$ est un anneau
noethérien, il en est de même de $S_0^\natural=S$
(\hyperref[II.2.1.5-fr]{2.1.5}). Inversement, si $S$ est une $S_0$-algèbre

